    \documentclass[conference]{IEEEtran}
    \IEEEoverridecommandlockouts
    % The preceding line is only needed to identify funding in the first footnote. If that is unneeded, please comment it out.
    \usepackage{cite}
    \usepackage{amsmath,amssymb,amsfonts}
    \usepackage{algorithmic}
    \usepackage{graphicx}
    \usepackage{textcomp}
    \usepackage{xcolor}
    \usepackage{placeins}
    \usepackage{float}
    \usepackage{url}
    \renewcommand{\abstractname}{Resumen}
    \renewcommand{\refname}{Referencias}
    \renewcommand{\IEEEkeywordsname}{Palabras claves}
    \def\BibTeX{{\rm B\kern-.05em{\sc i\kern-.025em b}\kern-.08em
        T\kern-.1667em\lower.7ex\hbox{E}\kern-.125emX}}
    \begin{document}

    \title{Detección de ataques de red multiclase sobre UNSW-NB15: análisis exploratorio y evaluación de modelos de Machine Learning}

    \author{\IEEEauthorblockN{1\textsuperscript{st} Fabricio Alonso Lanche Pacsi}
    \IEEEauthorblockA{\textit{dept. Computación} \\
    \textit{Ciencias de la Computación}\\
    Lima, Perú \\
    fabricio.lanche@utec.edu.pe}
    \and
    \IEEEauthorblockN{2\textsuperscript{nd} Jhogan Haldo Pachacutec Aguilar}
    \IEEEauthorblockA{\textit{dept. Computación} \\
    \textit{Ciencias de la Computación}\\
    Lima, Perú\\
    jhogan.pachacutec@utec.edu.pe}
    \and
    \IEEEauthorblockN{3\textsuperscript{rd} Anyeli Azumi Tamara Ureta}
    \IEEEauthorblockA{\textit{dept. Computación} \\
    \textit{Ciencias de la Computación}\\
    Lima, Perú \\
    anyeli.tamara@utec.edu.pe}
    \and
    \IEEEauthorblockN{4\textsuperscript{th} Paulo Isael Miranda Barrientos}
    \IEEEauthorblockA{\textit{dept. Computación} \\
    \textit{Ciencias de la Computación}\\
    Lima, Perú \\
    paulo.miranda@utec.edu.pe}
    \and
    \IEEEauthorblockN{5\textsuperscript{th} Christopher Renato Perez Torres}
    \IEEEauthorblockA{\textit{dept. Computación} \\
    \textit{Ciencias de Datos}\\
    Lima, Perú \\
    christopher.perez@utec.edu.pe}
    }

    \maketitle

    \begin{abstract}
    Este proyecto analiza el dataset UNSW-NB15 para desarrollar y evaluar modelos de Machine Learning capaces de clasificar tráfico de red en distintas categorías de ataques, considerando el desbalance de clases, la calidad de los datos y las características más relevantes para la detección de intrusiones.
    \end{abstract}

    \begin{IEEEkeywords}
    UNSW-NB15, clasificación, ciberseguridad, EDA.
    \end{IEEEkeywords}

    \section{Introducción y formulación del problema}

    \subsection{Motivación y relevancia del problema}
La detección de intrusiones en redes constituye un problema relevante 
en ciberseguridad, debido a la frecuencia y diversidad de los ataques 
que afectan actualmente a organizaciones y sistemas informáticos. 
Según el \textit{ENISA Threat Landscape 2025} \cite{enisa2025}, 
los ataques DDoS representaron aproximadamente el 77\,\% de los 
incidentes reportados durante el periodo analizado, mientras que el 
ransomware fue identificado como una de las amenazas de mayor impacto. 
Asimismo, el phishing estuvo involucrado en el 60\,\% de los puntos 
iniciales de intrusión, seguido por la explotación de vulnerabilidades 
con un 21,3\,\%.

Estas cifras muestran que las amenazas no se presentan de una única forma ni tienen el mismo comportamiento o impacto. Por ello, identificar solamente si un flujo de red es malicioso o benigno puede resultar insuficiente para caracterizar adecuadamente una amenaza. En este proyecto se plantea un problema de clasificación multiclase, orientado a identificar la categoría específica a la que pertenece cada flujo de tráfico. Distinguir entre diferentes tipos de comportamiento malicioso puede proporcionar información más detallada sobre un evento de seguridad y facilitar su análisis y la elección de medidas de respuesta adecuadas.

Para abordar este problema se utilizará el dataset \textbf{UNSW-NB15}, compuesto por registros de tráfico de red representados mediante variables que describen las características y el comportamiento de cada conexión. La variable objetivo del proyecto será \texttt{attack\_cat}, la cual identifica la categoría a la que pertenece cada registro, incluyendo distintos tipos de tráfico malicioso y una categoría \textit{Benign} para el tráfico no malicioso. A partir de estas variables, se buscará construir un modelo de clasificación multiclase capaz de identificar la categoría correspondiente a cada flujo de red.
\subsection{Introducción y formulación del problema}
Los sistemas de detección de intrusiones buscan identificar comportamientos anómalos o maliciosos dentro del tráfico de una red. Sin embargo, una clasificación limitada a únicamente determinar si un flujo es benigno o malicioso no permite conocer la naturaleza específica de la amenaza detectada. Distintos tipos de ataques pueden presentar patrones de comunicación diferentes y requerir medidas de respuesta distintas, por lo que identificar correctamente su categoría puede aportar información más útil para el análisis de un incidente de seguridad.
La relevancia de este problema radica en que una identificación más específica del tráfico malicioso puede contribuir al desarrollo de sistemas de detección de intrusiones con mayor capacidad de análisis. Reconocer el tipo de ataque puede facilitar la priorización de alertas, el estudio de patrones de comportamiento y la selección de medidas de respuesta acordes con la amenaza identificada. Desde la perspectiva de Machine Learning, además, el problema plantea el desafío de aprender patrones que permitan generalizar la clasificación a nuevos flujos de red no observados durante el entrenamiento.


\subsection{Objetivo y tarea de Machine Learning}

Este problema se abordará mediante un enfoque de aprendizaje supervisado, dado que el conjunto de datos contiene la ``verdad fundamental'' para cada uno de los registros históricos evaluados, la cual se encuentra representada explícitamente mediante la etiqueta atack\_cat (variable objetivo). Al contar con este emparejamiento, los modelos se optimizarán para aprender los patrones topológicos y de flujo de la red, tales como la duración de la conexión, el conteo de paquetes y los protocolos, y asociarlos con las etiquetas conocidas. Por este motivo, se descarta un enfoque no supervisado, como el \textit{clustering}, ya que el propósito no es descubrir estructuras o anomalías desconocidas, sino aprender a partir de ejemplos previamente categorizados.

Para abordar esta tarea, se empleará una Regresión Logística Multinomial como modelo \textit{baseline}, debido a su estructura lineal y a su utilidad como referencia para evaluar modelos de mayor complejidad. Como modelos de comparación, se utilizarán una Máquina de Vectores de Soporte (SVM) y un Random Forest, lo que permitirá contrastar un modelo lineal con un clasificador basado en maximización de márgenes y con un método de ensamble basado en múltiples árboles. Esta selección cuenta con precedentes en estudios de clasificación sobre UNSW-NB15, donde se han evaluado conjuntamente modelos de Regresión Logística y métodos de clasificación basados en árboles mediante métricas de clasificación y matrices de confusión \cite{more2024enhanced}.

El entrenamiento seguirá un esquema de partición entrenamiento-validación-prueba, destinando el 75\% de los datos al entrenamiento, el 15\% a la validación y el 10\% restante a la evaluación final.

El objetivo principal del proyecto es predecir y clasificar el tráfico de red con el fin de identificar incidentes de ciberseguridad. Concretamente, la tarea de Machine Learning consiste en construir y entrenar un modelo de clasificación multiclase capaz de categorizar una conexión de red en uno de 10 estados posibles: tráfico normal (\texttt{Benign}) o una de las 9 familias de ataques cibernéticos  (\texttt{Generic}, \texttt{Exploits}, \texttt{Fuzzers}, \texttt{DoS}, \texttt{Reconnaissance}, \texttt{Analysis}, \texttt{Backdoors}, \texttt{Shellcode} y \texttt{Worms}).


\subsection{Preguntas de investigación}

\textbf{Pregunta de investigación principal:}

¿En qué medida las características de los flujos de tráfico de red del dataset UNSW-NB15 permiten clasificar correctamente cada registro como una de las distintas categorías de ataque mediante técnicas de Machine Learning?

\textbf{Preguntas de investigación específicas:}

\begin{itemize}
    \item ¿Qué características del tráfico de red presentan mayor relación con la categoría de ataque identificada por \texttt{attack\_cat}?
    
    \item ¿Existen categorías de tráfico que presenten patrones similares y, por lo tanto, sean más difíciles de distinguir entre sí?
    
    \item ¿Cómo influye la distribución de las clases en la capacidad de identificar correctamente las distintas categorías de tráfico?
\end{itemize}
\FloatBarrier

\section{Datos y Análisis Exploratorio de Datos (EDA)}

Esta sección presenta el análisis exploratorio del proyecto. Los scripts y el código necesarios para reproducir el análisis se encuentran disponibles en el repositorio del proyecto \cite{repo}.

\subsection{Descripción del dataset}

\FloatBarrier

El dataset utilizado para este trabajo es UNSW-NB15, generado por 
Moustafa y Slay \cite{moustafa2015unsw} a partir de tráfico de red 
capturado en el Cyber Range Lab de la ACCS. Un aspecto relevante 
para su interpretación es que las observaciones no provienen de una 
sola sesión de captura homogénea, sino de una combinación deliberada 
de actividad normal moderna real y comportamientos de ataque 
sintéticos generados con la herramienta IXIA PerfectStorm 
\cite{moustafa2015unsw}. Los propios autores señalan que este diseño 
buscó superar las limitaciones de datasets más antiguos como KDD99 y 
NSL-KDD, cuyo tráfico normal y patrones de ataque ya no reflejan el 
comportamiento de redes contemporáneas \cite{moustafa2016evaluation}.

El dataset completo está compuesto por 2{,}540{,}044 registros 
distribuidos originalmente en cuatro archivos CSV 
\cite{moustafa2015unsw}. Los autores también publicaron una partición 
reducida en conjuntos de entrenamiento y prueba (175{,}341 y 82{,}332 
registros, respectivamente, con una proporción aproximada de 60:40) 
\cite{moustafa2016evaluation}; sin embargo, en este proyecto se 
utiliza la versión completa del dataset, disponible públicamente en 
Kaggle \cite{kaggle_unsw}, con el fin de contar con una mayor cantidad 
de observaciones para el entrenamiento y una representación más 
amplia de las clases minoritarias. 

Cada registro representa un flujo de red bidireccional (conexión 
cliente-servidor o servidor-cliente) caracterizado mediante variables 
extraídas con las herramientas Argus y Bro-IDS, junto con doce 
algoritmos adicionales desarrollados por los propios autores para 
generar variables derivadas \cite{moustafa2015unsw}. El resultado de 
este proceso es un total de 49 columnas por registro, entre variables 
predictoras y variables de etiqueta.

\FloatBarrier

\subsection{Variables y variable objetivo}

Las 49 columnas del dataset se organizan conceptualmente en cinco 
grupos, de acuerdo con la taxonomía definida por los propios autores 
\cite{moustafa2016evaluation}: variables de flujo, variables básicas, 
variables de contenido, variables de tiempo y variables adicionales 
generadas. La Tabla \ref{tab:grupos_variables} resume esta 
organización.

\FloatBarrier

\begin{table}[H]
\centering
\caption{Grupos de variables del dataset UNSW-NB15 (Moustafa \& 
Slay, 2016).}
\label{tab:grupos_variables}
\footnotesize
\begin{tabular}{@{}p{1.6cm}p{0.7cm}p{4.3cm}@{}}
\hline
\textbf{Grupo} & \textbf{N.} & \textbf{Ejemplos} \\
\hline
Flujo & 5 & srcip, dstip, sport, proto \\
Básicas & 13 & state, dur, sbytes, sttl \\
Contenido & 8 & swin, smeansz, trans\_depth \\
Tiempo & 9 & sjit, sintpkt, tcprtt \\
Adicionales & 12 & ct\_state\_ttl, ct\_srv\_src \\
\hline
Etiquetas & 2 & attack\_cat, label \\
\hline
\textbf{Total} & \textbf{49} & \\
\hline
\end{tabular}
\end{table}

Así pues, cada grupo captura un nivel distinto de información sobre 
la conexión de red:

\begin{itemize}
    \item Las variables de \textbf{flujo} identifican unívocamente 
    la conexión (direcciones IP, puertos y protocolo de transporte), 
    sin describir por sí mismas su comportamiento.
    \item Las variables \textbf{básicas} resumen las propiedades 
    generales del flujo, como su duración, el volumen de bytes 
    intercambiados en cada dirección y el tiempo de vida (TTL) de 
    los paquetes.
    \item Las variables de \textbf{contenido} profundizan en las 
    cabeceras TCP y en la estructura de las transacciones HTTP, 
    aportando detalle sobre el contenido transmitido dentro de la 
    conexión.
    \item Las variables de \textbf{tiempo} describen la dinámica 
    temporal del flujo, incluyendo el jitter y los tiempos de 
    establecimiento de la conexión TCP (handshake).
    \item Las variables \textbf{adicionales generadas} no provienen 
    directamente de un único paquete o conexión, sino que se 
    construyen a partir de ventanas de hasta 100 registros 
    consecutivos, contabilizando la frecuencia con la que se repiten 
    ciertas combinaciones de IP, puerto o servicio; este diseño 
    busca capturar patrones de comportamiento repetitivo 
    característicos de ciertos tipos de ataque, como los escaneos o 
    las conexiones masivas de corta duración.
\end{itemize}

Respecto al tipo de dato, las variables del dataset se dividen 
principalmente en dos tipos: numéricas (de tipo entero o de punto 
flotante, representando magnitudes como duración, bytes o tiempos 
entre paquetes) y nominales o categóricas (de tipo cadena de texto, 
como \texttt{proto}, \texttt{state} y \texttt{service}) 
\cite{moustafa2016evaluation}. Los propios autores destacan que esta 
falta de un formato estándar entre variables numéricas y nominales 
constituye un desafío recurrente para las técnicas de detección de 
intrusiones basadas en aprendizaje automático 
\cite{moustafa2016evaluation}. Este dataset esta compuesto por 41 variables 
numéricas y un subconjunto de variables nominales de alta 
cardinalidad (\texttt{srcip}, \texttt{dstip}, \texttt{sport}, 
\texttt{dsport}) además de las de baja cardinalidad mencionadas 
(\texttt{proto}, \texttt{state}, \texttt{service}).

Para el etiquetado del dataset, los autores proporcionan dos variables de clase \cite{moustafa2016evaluation}: \texttt{label}, una codificación binaria donde 0 indica tráfico normal y 1 indica  tráfico de ataque; y \texttt{attack\_cat}, que identifica la categoría específica del ataque entre nueve posibilidades:

\begin{itemize}
    \item \textbf{Fuzzers}: envío masivo de datos aleatorios o  malformados a un programa o servicio con el fin de provocar fallos y descubrir vulnerabilidades.
    \item \textbf{Analysis}: intrusiones diversas dirigidas a aplicaciones web, que incluyen escaneo de puertos, spam y  scripts embebidos en archivos HTML.
    \item \textbf{Backdoors}: técnicas para eludir mecanismos de autenticación y obtener acceso remoto no autorizado y sigiloso a un dispositivo.
    \item \textbf{DoS} (Denegación de Servicio): ataques que buscan saturar los recursos de memoria o procesamiento de un sistema para impedir el acceso de usuarios legítimos.
    \item \textbf{Exploits}: secuencias de instrucciones que aprovechan una falla o vulnerabilidad específica de un sistema o red para provocar un comportamiento no previsto.
    \item \textbf{Generic}: técnica que ataca cualquier esquema de cifrado por bloques mediante colisiones basadas en funciones hash, independientemente de la configuración del cifrado.
    \item \textbf{Reconnaissance}: actividades de sondeo (probing) orientadas a recolectar información sobre una red para evadir sus controles de seguridad.
    \item \textbf{Shellcode}: ataque en el que se introduce un fragmento de código malicioso que abre una shell para tomar control del equipo comprometido.
    \item \textbf{Worms}: malware autorreplicable que se propaga de forma autónoma a través de la red explotando vulnerabilidades, sin requerir interacción del usuario.
\end{itemize}
De este modo, se adopta \texttt{attack\_cat} como la variable objetivo del proyecto, planteando un problema de clasificación multiclase. Bajo esta formulación, la categoría \textit{Benign} representa el flujo de red normal o legítimo, mientras que las nueve clases descritas anteriormente permiten distinguir la tipología específica de cada amenaza, superando la limitación de un esquema binario simple.
\FloatBarrier

\subsection{Valores faltantes}

\FloatBarrier

El conjunto de datos presenta datos faltantes en 6 de sus 49 variables, afectando al 57.20\% de las observaciones. La Tabla~\ref{tab:faltantes} desglosa la cantidad y el origen de los valores faltantes por variable.

\begin{table}[h!]
\centering
\caption{Valores faltantes por columna.}
\label{tab:faltantes}
\begin{tabular}{lrrl}
\hline
\textbf{Columna} & \textbf{Faltantes} & \textbf{\%} & \textbf{Origen} \\
\hline
\texttt{is\_ftp\_login} & 1,430,065 & 56.30 & Protocolo FTP \\
\texttt{ct\_ftp\_cmd}    & 1,429,879 & 56.29 & Protocolo FTP \\
\texttt{ct\_flw\_http\_mthd} & 1,348,145 & 53.08 & Protocolo HTTP \\
\texttt{dsport}          & 304       & 0.01  & Token corrupto \\abordará
\texttt{sport}           & 8         & 0.00  & Token corrupto \\
\texttt{state}           & 8         & 0.00  & Token corrupto \\
\hline
\end{tabular}
\end{table}

Los faltantes responden a dos orígenes distintos. Las tres variables con mayor proporción de datos faltantes están directamente asociadas al protocolo de aplicación. Cada flujo en el dataset se etiqueta con su protocolo en la variable \texttt{service} (HTTP, FTP, DNS, etc.), y las variables \texttt{is\_ftp\_login} y \texttt{ct\_ftp\_cmd} solo tienen sentido semántico en sesiones FTP, mientras que \texttt{ct\_flw\_http\_mthd} solo aplica a sesiones HTTP. Estos faltantes son valores que no existen por el propio diseño de los datos, no suponen ser problemas de captura.

Los nulos restantes en \texttt{sport} (8), \texttt{dsport} (304) y \texttt{state} (8) representan menos del 0.02\% del total y provienen de tokens corruptos en la fuente original (literales hexadecimales ilegibles, guiones y espacios). Fueron marcados intencionalmente como \texttt{NaN} durante la carga de datos.

\FloatBarrier

\subsection{Distribución de variables y clases}

\FloatBarrier

La variable objetivo \texttt{attack\_cat} define una tarea de clasificación multiclase con 10 categorías: las 9 familias de ataque de UNSW-NB15 más la clase artificial \texttt{Benign}, que agrupa el tráfico de red normal sin etiqueta de ataque específica (identificado con \texttt{Label} = 0).

La Tabla~\ref{tab:descriptivos} resume los estadísticos descriptivos de una selección de variables representativas por grupo funcional, de acuerdo con la Tabla~\ref{tab:grupos_variables}, e incluye tanto variables con distribuciones extremas como variables de contraste. Muchas exhiben asimetrías pronunciadas, con medias sustancialmente superiores a sus medianas, por ejemplo \texttt{Sjit} (1{,}589 frente a 19.12) o \texttt{Sintpkt} (193.32 frente a 0.47); otras presentan alta varianza en magnitudes muy dispares, como \texttt{sbytes} y \texttt{sport}. Estas distribuciones sesgadas resultan típicas del tráfico de red y confirman la necesidad de normalización durante el preprocesamiento. El detalle completo de las 41 variables numéricas está disponible en el repositorio del proyecto (\texttt{output/descriptivos.csv}).

\begin{table}[h!]
\centering
\caption{Estadísticos descriptivos de variables numéricas seleccionadas por grupo funcional.}
\label{tab:descriptivos}
\resizebox{\columnwidth}{!}{%
\begin{tabular}{llrrrrr}
\hline
\textbf{Grupo} & \textbf{Variable} & \textbf{Media} & \textbf{Mediana} & \textbf{Varianza} & \textbf{Mín} & \textbf{Máx} \\
\hline
Básicas & \texttt{dur} & 0.66 & 0.02 & 193.90 & 0.00 & 8,787 \\
Básicas & \texttt{sbytes} & 4,340 & 1,470 & 3.182e+09 & 0.00 & 1.436e+07 \\
Básicas & \texttt{sttl} & 62.78 & 31.00 & 5,569 & 0.00 & 255.00 \\
Contenido & \texttt{trans\_depth} & 0.08 & 0.00 & 0.12 & 0.00 & 172.00 \\
Tiempo & \texttt{Sjit} & 1,589 & 19.12 & 2.860e+08 & 0.00 & 1,483,831 \\
Tiempo & \texttt{Sintpkt} & 193.32 & 0.47 & 7,723,748 & 0.00 & 84,371 \\
Tiempo & \texttt{tcprtt} & 0.0062 & 0.0006 & 0.0021 & 0.00 & 10.04 \\
Adicionales & \texttt{ct\_state\_ttl} & 0.26 & 0.00 & 0.47 & 0.00 & 6.00 \\
Adicionales & \texttt{ct\_srv\_src} & 9.21 & 5.00 & 117.44 & 1.00 & 67.00 \\
\hline
\end{tabular}}
\end{table}

La Fig.~\ref{fig:boxplot} complementa los estadísticos de la
Tabla~\ref{tab:descriptivos} mostrando la dispersión de las variables
seleccionadas en escala logarítmica, lo que permite apreciar su asimetría y
rango dinámico. La detección de outliers y su cuantificación se abordan en la subsección siguiente.


\begin{figure}[h!]
\centering
\includegraphics[width=\columnwidth]{output/fig02_boxplot_outliers.png}
\caption{Boxplots en escala logarítmica de las variables numéricas seleccionadas en la Tabla~\ref{tab:descriptivos}.}
\label{fig:boxplot}
\end{figure}

\begin{figure}[h!]
\centering
\includegraphics[width=\columnwidth]{output/fig01_distribucion_clases.png}
\caption{Distribución de la variable objetivo: las 10 clases (izquierda) y zoom sobre las 9 familias de ataque (derecha)}
\label{fig:clases}
\end{figure}

Por otro lado, La Fig.~\ref{fig:clases} (izquierda) muestra que el 87.35\% de los registros corresponden a \texttt{Benign}, mientras que los ataques representan apenas el 12.65\% del total. El panel derecho, que excluye a \texttt{Benign} para apreciar el detalle de las clases minoritarias, revela que el tráfico malicioso se concentra principalmente en \texttt{Generic} (215\,481 registros, 8.48\,\%), seguido de \texttt{Exploits}, \texttt{Fuzzers}, \texttt{DoS} y \texttt{Reconnaissance}; en el extremo opuesto, \texttt{Worms} cuenta con solo 174 registros. Con esta distribución, los modelos de clasificación tenderán naturalmente a favorecer la predicción de la clase mayoritaria, requiriendo técnicas específicas de balanceo o evaluación que penalicen el error en clases minoritarias durante el entrenamiento.

\subsection{Problemas de calidad y outliers}

\FloatBarrier

Más allá de los faltantes estructurales, el dataset presenta otros problemas de calidad que requieren atención. Se identificaron 480,633 filas duplicadas (18.92\% del total), correspondientes a registros con valores idénticos en todas las características y el mismo target. Estos duplicados exactos no aportan información nueva al modelo y son candidatos naturales para eliminación en el preprocesamiento. No se detectaron columnas constantes.

La Tabla~\ref{tab:cardinalidad} presenta la cardinalidad de las variables nominales. Como se observa, las direcciones IP \texttt{srcip} y \texttt{dstip} presentan únicamente 43 y 47 valores únicos respectivamente, una diversidad inusualmente baja que sugiere un escenario de captura con pocos hosts activos. Los puertos, por el contrario, exhiben una alta cardinalidad: \texttt{sport} y \texttt{dsport} registran aproximadamente 65 mil valores únicos cada uno (2.5\% del total), reflejando la variabilidad esperada en puertos efímeros TCP/UDP. El resto de variables nominales (\texttt{proto}, \texttt{state} y \texttt{service}) se limita a pocas categorías distintas.

\begin{table}[h!]
\centering
\caption{Cardinalidad de las variables nominales.}
\label{tab:cardinalidad}
\begin{tabular}{lrr}
\hline
\textbf{Variable} & \textbf{Valores únicos} & \textbf{\% del total} \\
\hline
\texttt{srcip}   & 43      & 0.002 \\
\texttt{dstip}   & 47      & 0.002 \\
\texttt{sport}   & 64,598  & 2.543 \\
\texttt{dsport}  & 64,627  & 2.544 \\
\texttt{proto}   & 135     & 0.005 \\
\texttt{state}   & 16      & 0.001 \\
\texttt{service} & 13      & 0.001 \\
\hline
\end{tabular}
\end{table}

Para detectar valores atípicos en variables numéricas se aplicó el criterio del rango intercuartílico (IQR), clasificando como outliers todo valor fuera del intervalo $[Q_1 - 1.5 \times \mathrm{IQR}, Q_3 + 1.5 \times \mathrm{IQR}]$. La Tabla~\ref{tab:outliers} cuantifica el porcentaje de outliers y la proporción de ceros para cada variable numérica, ordenadas por prevalencia de atípicos.

Una fracción elevada de variables numéricas presenta conteos significativos de outliers, pero estos reflejan en general la forma natural de las distribuciones más que corrupción de datos. Los outliers se concentran en variables con colas largas y asimétricas, tales como \texttt{Sload} (22.17\% outliers, 0.20\% ceros), \texttt{Sjit} (18.63\% outliers, 37.93\% ceros) y \texttt{dur} (17.77\% outliers, 0.33\% ceros), que contrastan con variables acotadas como los campos TTL (\texttt{sttl}, \texttt{dttl}) y la mayoría de contadores \texttt{ct\_*}, que no presentan outliers bajo este criterio. Por otro lado, variables como \texttt{ct\_ftp\_cmd} (96.07\% ceros), \texttt{res\_bdy\_len}  (92,65\%) y \texttt{trans\_depth} (91.92\%) tienen una proporción mayoritaria de ceros, reforzando la interpretación de que estas métricas no son relevantes para todas las conexiones.

\begin{table}[h!]
\centering
\caption{Porcentaje de outliers (criterio IQR) y proporción de ceros en las variables numéricas.}
\label{tab:outliers}
\footnotesize
\begin{tabular}{lrr}
\hline
\textbf{Variable} & \textbf{\% outliers} & \textbf{\% ceros} \\
\hline
\texttt{Sload}             & 22.17 & 0.20  \\
\texttt{ct\_src\_dport\_ltm} & 21.01 & 0.00  \\
\texttt{Sjit}              & 18.63 & 37.93 \\
\texttt{Dintpkt}           & 17.86 & 19.71 \\
\texttt{dur}               & 17.77 & 0.33  \\
\texttt{Sintpkt}           & 17.56 & 0.06  \\
\texttt{ct\_dst\_src\_ltm} & 16.36 & 0.00  \\
\texttt{Djit}              & 16.30 & 35.43 \\
\texttt{ct\_dst\_ltm}      & 14.28 & 0.00  \\
\texttt{dsport}            & 13.88 & 2.13  \\
\texttt{Dload}             & 13.57 & 19.71 \\
\texttt{ct\_src\_ltm}      & 13.26 & 0.00  \\
\texttt{ct\_srv\_src}      & 12.36 & 0.00  \\
\texttt{ct\_srv\_dst}      & 12.23 & 0.00  \\
\texttt{dbytes}            & 11.89 & 19.69 \\
\texttt{sbytes}            & 11.34 & 0.00  \\
\texttt{sloss}             & 8.03  & 41.23 \\
\texttt{smeansz}           & 6.66  & 0.00  \\
\texttt{Dpkts}             & 5.92  & 19.69 \\
\texttt{Spkts}             & 5.62  & 0.00  \\
\texttt{dloss}             & 5.08  & 41.36 \\
\texttt{ackdat}            & 4.93  & 41.55 \\
\texttt{tcprtt}            & 4.42  & 41.54 \\
\texttt{synack}            & 4.38  & 41.54 \\
\texttt{dmeansz}           & 1.01  & 19.69 \\
\texttt{dttl}              & 0.00  & 19.92 \\
\texttt{sport}             & 0.00  & 2.13  \\
\texttt{sttl}              & 0.00  & 0.41  \\
\texttt{dwin}              & 0.00  & 41.28 \\
\texttt{swin}              & 0.00  & 41.14 \\
\texttt{stcpb}             & 0.00  & 41.25 \\
\texttt{trans\_depth}      & 0.00  & 91.92 \\
\texttt{Stime}             & 0.00  & 0.00  \\
\texttt{res\_bdy\_len}     & 0.00  & 92.65 \\
\texttt{dtcpb}             & 0.00  & 41.27 \\
\texttt{Ltime}             & 0.00  & 0.00  \\
\texttt{ct\_flw\_http\_mthd} & 0.00 & 82.79 \\
\texttt{ct\_state\_ttl}    & 0.00  & 85.75 \\
\texttt{ct\_ftp\_cmd}      & 0.00  & 96.07 \\
\texttt{ct\_dst\_sport\_ltm} & 0.00 & 0.00  \\
\hline
\end{tabular}
\end{table}

\FloatBarrier

\subsection{Relaciones entre variables (correlaciones)}

\FloatBarrier

El análisis de relaciones entre variables permite identificar patrones de asociación con el target y detectar redundancia en el espacio de características. La Fig.~\ref{fig:correlacion} muestra las correlaciones de Pearson entre las variables numéricas y el target binario. La variable \texttt{sttl} (0.90) destaca por su correlación excepcional con el target, mientras que el resto de las variables incluidas en esa figura presenta correlaciones débiles o moderadas. Entre estas, \texttt{Dload} muestra la mayor correlación negativa (--0.22), sugiriendo que variables de volumen de datos descendente son menos indicativas de ataque que otras características. Los contadores de ventana \texttt{ct\_*}, incluido \texttt{ct\_state\_ttl}, no se incorporan en esa figura de correlación general; su relación con el target se analiza por separado en la Fig.~\ref{fig:ct_corr}.

\begin{figure}[h!]
\centering
\includegraphics[width=\columnwidth]{output/fig03_correlacion_features.png}
\caption{Correlación de Pearson entre las variables numéricas y el target binario.}
\label{fig:correlacion}
\end{figure}

La Fig.~\ref{fig:hist} complementa este análisis al mostrar cómo se distribuyen las variables clave según la clase. Observando las separaciones entre la curva de tráfico benigno y la de ataque, variables como \texttt{sttl} y \texttt{ct\_state\_ttl} exhiben diferenciación clara entre poblaciones, mientras que otras como \texttt{sport} mantienen distribuciones prácticamente idénticas.

\begin{figure}[h!]
\centering
\includegraphics[width=\columnwidth]{output/fig04_histogramas_top9.png}
\caption{Distribución de variables numéricas clave según la clase.}
\label{fig:hist}
\end{figure}

Más allá de la correlación con el target, el espacio de características presenta alta redundancia entre predictores. Pares como \texttt{Stime} y \texttt{Ltime} (1.00), \texttt{swin} y \texttt{dwin} (0.997), y \texttt{sbytes} con \texttt{sloss} (0.96) transmiten información duplicada, requiriendo selección cuidadosa durante el modelado. Los contadores de ventana \texttt{ct\_*} presentan las correlaciones más altas con el target, encabezados por \texttt{ct\_state\_ttl} (0.874) y \texttt{ct\_dst\_src\_ltm} (0.440), como se detalla en la Fig.~\ref{fig:ct_corr}. 
\begin{figure}[h!]
\centering
\includegraphics[width=\columnwidth]{output/fig05_ct_vs_target.png}
\caption{Correlación de las variables de ventana \texttt{ct\_*} con el target.}
\label{fig:ct_corr}
\end{figure}

En el contexto multiclase, la Tabla~\ref{tab:eta} reporta la razón de correlación $\eta^{2}$, una medida estándar de asociación entre una variable numérica y una variable categórica que no supone una relación lineal (a diferencia de la correlación de Pearson). Cuantifica la proporción de la varianza de la variable numérica que es explicada por las diferencias entre los grupos de la variable categórica, y se define como
\begin{equation}
    \eta^{2} = \frac{\mathrm{SSB}}{\mathrm{SST}} = \frac{\sum_{g} n_{g}(\bar{y}_{g} - \bar{y})^{2}}{\sum_{i} (y_{i} - \bar{y})^{2}},
\end{equation}
donde $\bar{y}_{g}$ y $n_{g}$ son la media y el número de observaciones del grupo $g$ (cada clase de \texttt{attack\_cat}), y $\bar{y}$ es la media global; toma valores entre 0 y 1, y un valor alto indica que la variable numérica discrimina fuertemente entre las 10 clases. En este contexto los mismos predictores dominan, con \texttt{sttl} ($\eta^2 = 0.835$) y \texttt{ct\_state\_ttl} (0.777) nuevamente a la cabeza, seguidos por contadores \texttt{ct\_*} en el rango 0.22--0.32. El resto de variables presenta $\eta^2$ bajo ($ \leq 0.14$), indicando que la discriminación entre las 10 clases de ataque depende de un subconjunto reducido de características, mientras que la mayoría de variables aporta poco en el contexto multiclase.

\begin{table}[h!]
\centering
\caption{Razón de correlación $\eta^2$ entre las variables numéricas y \texttt{attack\_cat} (10 clases).}
\label{tab:eta}
\footnotesize
\begin{tabular}{lrlr}
\hline
\textbf{Variable} & \textbf{$\eta^2$} & \textbf{Variable} & \textbf{$\eta^2$} \\
\hline
\texttt{sttl}                & 0.835 & \texttt{dsport}              & 0.049 \\
\texttt{ct\_state\_ttl}      & 0.777 & \texttt{Dload}               & 0.048 \\
\texttt{ct\_dst\_src\_ltm}   & 0.324 & \texttt{Sload}               & 0.041 \\
\texttt{dttl}                & 0.307 & \texttt{sport}               & 0.030 \\
\texttt{ct\_dst\_sport\_ltm} & 0.301 & \texttt{smeansz}             & 0.024 \\
\texttt{ct\_srv\_dst}        & 0.290 & \texttt{Spkts}               & 0.016 \\
\texttt{ct\_srv\_src}        & 0.283 & \texttt{Sjit}                & 0.014 \\
\texttt{ct\_src\_dport\_ltm} & 0.274 & \texttt{Dpkts}               & 0.014 \\
\texttt{ct\_dst\_ltm}        & 0.238 & \texttt{ct\_flw\_http\_mthd} & 0.012 \\
\texttt{ct\_src\_ltm}        & 0.226 & \texttt{dloss}               & 0.010 \\
\texttt{swin}                & 0.133 & \texttt{dur}                 & 0.009 \\
\texttt{dwin}                & 0.132 & \texttt{trans\_depth}        & 0.008 \\
\texttt{tcprtt}              & 0.105 & \texttt{sloss}               & 0.006 \\
\texttt{ackdat}              & 0.092 & \texttt{dbytes}              & 0.006 \\
\texttt{synack}              & 0.090 & \texttt{sbytes}              & 0.004 \\
\texttt{dmeansz}             & 0.079 & \texttt{ct\_ftp\_cmd}        & 0.004 \\
\texttt{Stime}               & 0.078 & \texttt{Djit}                & 0.003 \\
\texttt{Ltime}               & 0.078 & \texttt{res\_bdy\_len}       & 0.001 \\
\texttt{dtcpb}               & 0.073 & \texttt{Sintpkt}             & 0.001 \\
\texttt{stcpb}               & 0.073 & \texttt{Dintpkt}             & 0.001 \\
\hline
\end{tabular}
\end{table}

\FloatBarrier

\subsection{Sesgo y data leakage}

A partir del análisis exploratorio, se evalúan los fenómenos de sesgo y fuga de información. El \emph{sesgo} comprende cualquier propiedad de los datos que desvíe al modelo del comportamiento general del tráfico operativo \cite{mehrabi2021survey}. Por otro lado, la \emph{fuga de información} (\emph{data leakage}) consiste en la inclusión de características no disponibles en el momento real de la predicción, introduciendo canales no deseados hacia la variable objetivo \cite{kaufman2012leakage}. Mientras el sesgo degrada la capacidad predictiva del modelo, la fuga genera una sobreestimación ilusoria de su desempeño.
\subsubsection{Identificación y evaluación de Sesgos}

La distribución de la variable objetivo introduce el primer sesgo, y el más evidente. El 87.35\% de los registros corresponde a tráfico Benign, mientras que los ataques representan solo el 12.65\% y se concentran principalmente en la familia Generic (Fig.~\ref{fig:clases}). La situación se agrava en el interior de las clases de ataque: la familia Worms cuenta con apenas 174 registros. Un clasificador multiclase entrenado sobre estas proporciones tenderá a favorecer la clase mayoritaria y a subrepresentar los ataques poco frecuentes. La exactitud global ocultaría ese comportamiento, por lo que la evaluación por clase resultará indispensable en la siguiente etapa.
El segundo sesgo se constituye por un sesgo de sobrerrepresentacion, donde existe un 18.92\% de filas duplicadas (480\,633) . Los flujos repetidos no aportan información nueva, pero cuentan con mayor peso durante el entrenamiento y, al reforzar los patrones mayoritarios, amplifican el desbalance descrito al comienzo. Su eliminación queda prevista en el preprocesamiento.

Adicionalmente, se identifica un \emph{sesgo de representatividad} derivado de la captura en un entorno controlado. Las direcciones IP presentan una cardinalidad inusualmente baja (43 orígenes y 47 destinos, Tabla~\ref{tab:cardinalidad}), lo que evidencia un escenario con pocos anfitriones activos. Este sesgo introduce el \emph{riesgo de memorización} y asociación espuria entre la identidad de un host y una familia de ataque específica: el clasificador arriesga memorizar direcciones IP en lugar de aprender los patrones de tráfico subyacentes, imposibilitando su generalización ante nuevas topologías de red.

Por otra parte, existe un sesgo a causa del dominio de valores extremos dentro de las variables numéricas. Según la Tabla~\ref{tab:descriptivos}, variables como \texttt{Sload} exhiben asimetrías severas (media $3.7\times10^{7}$ frente a una mediana de $5.9\times10^{5}$) y una proporción elevada de observaciones clasificadas como atípicas bajo el criterio IQR (22.2\%); lo mismo ocurre en \texttt{Sjit} (18.6\%) y \texttt{dur} (17.8\%) (Tabla~\ref{tab:outliers}). Estos valores extremos constituyen la forma natural de las distribuciones y no indican corrupción de los datos, pero sin una transformación previa dominan el aprendizaje, que pasaría a depender de la magnitud del flujo en lugar de su comportamiento. Se prevé una normalización logarítmica en la etapa de preprocesamiento.

Por último, el espacio de características contiene predictores redundantes. Los pares \texttt{Stime}--\texttt{Ltime} (correlación 1.00), \texttt{swin}--\texttt{dwin} (0.997) y \texttt{sbytes}--\texttt{sloss} (0.96) transmiten información duplicada. Además, el análisis multiclase mediante la razón de correlación $\eta^{2}$ indica que solo \texttt{sttl} (0.835) y \texttt{ct\_state\_ttl} (0.777) discriminan con fuerza entre las diez categorías, mientras que el resto obtiene valores inferiores o iguales a 0.14 (Tabla~\ref{tab:eta}). La redundancia sobrepondera los factores duplicados, y la gran cantidad de variables poco informativas añade ruido sin aportar capacidad de clasificación; ambos efectos se atenderán con selección de características.

\subsubsection{Data leakage}

La fuga de información se manifiesta en este dataset a través de tres fuentes, todas advertidas a partir del análisis exploratorio. La más directa tiene su origen en la propia construcción de la variable objetivo. La clase Benign se deriva de la etiqueta binaria \texttt{Label} (un flujo es Benign si \texttt{Label} $=$ 0), por lo que \texttt{attack\_cat} y \texttt{Label} codifican la misma información bajo dos formas: cualquier variable derivada de la etiqueta que ingresara al conjunto de características filtraría la respuesta esperada y el modelo aprendería a leer la salida en lugar de la señal de los datos. La separación entre variables de entrada y objetivo debe, por tanto, excluir explícitamente ambas columnas del modelado.


Una segunda fuente proviene de las características de ventana temporal de la familia \texttt{ct\_*} (por ejemplo, \texttt{ct\_srv\_src}, \texttt{ct\_dst\_ltm}, \texttt{ct\_src\_dport\_ltm}), que agregan conteos de actividad alrededor de cada conexión. En el momento real de la decisión, en un sistema de detección en línea, no toda esa actividad es aún conocida: si un conteo incorpora flujos posteriores a la conexión que se clasifica, el modelo aprende con información del futuro. La misma observación aplica a los campos de tiempo \texttt{Stime} y \texttt{Ltime}, que registran el instante de la captura y cuya inclusión como predictores haría depender al modelo del momento en que ocurrió el flujo y no de su contenido.

Estas dos últimas fuentes comparten una consecuencia práctica: una partición aleatoria de los datos no las controla. Si el entrenamiento contiene conexiones temporalmente posteriores a las de prueba, el modelo puede aprovechar patrones que en producción jamás estarían disponibles. Para un dataset de tráfico de red, la partición debe respetar el orden temporal, de modo que las conexiones de prueba correspondan siempre a momentos posteriores a las de entrenamiento; esta decisión se prevé para la etapa de modelado.

Finalmente, los registros duplicados constituyen una vía de fuga entre los conjuntos de partición. El 18.92\% de filas duplicadas (480\,633) son ejemplos idénticos en todas las características y con la misma clase: en un reparto aleatorio, copias de un mismo flujo pueden quedar a la vez en entrenamiento y en prueba, y el modelo resultaría evaluado sobre ejemplos que ya vio durante el ajuste. La eliminación de duplicados antes de particionar cierra esta vía, y refuerza además la corrección del sesgo de sobrerrepresentación descrito en la subsección anterior.

En conclusión, el sesgo y fuga de información, pese a tener orígenes distintos, convergen en una misma consecuencia: distorsionar la relación entre el rendimiento reportado y el rendimiento real del modelo, comprometiendo desde el origen la confiabilidad del clasificador multiclase.


\section{Siguientes pasos}

Con base en el diagnóstico del análisis exploratorio (EDA), la siguiente fase del proyecto se estructurará en tres etapas consecutivas orientadas al desarrollo y evaluación de los modelos de clasificación.

En primer lugar, se ejecutará el \textbf{preprocesamiento e ingeniería de datos}. Esta etapa abarcará la eliminación previa de duplicados, la exclusión explícita de variables propensas a \emph{data leakage} o asociación espuria (\texttt{Label}, \texttt{attack\_cat} e IPs), y la aplicación de transformaciones logarítmicas para estabilizar variables numéricas con sesgo extremo. Asimismo, se seleccionarán las características más relevantes para reducir la redundancia y se definirá una partición estrictamente temporal, integrada en un \emph{pipeline} único ajustado exclusivamente sobre el conjunto de entrenamiento.

En segundo lugar, se desarrollará el estudio teórico y la implementación de los modelos seleccionados, la regresión logística multinomial, la Máquina de Vectores de Soporte (SVM) y el Random Forest como modelos de comparación. Esta etapa permitirá establecer las características y supuestos principales de cada modelo en relación con la tarea de clasificación multiclase planteada.

Finalmente, se realizará el ajuste de los modelos como parte del proceso de aprendizaje supervisado. El entrenamiento se llevará a cabo respetando la partición definida y las medidas establecidas para evitar \emph{data leakage}. Los modelos serán posteriormente evaluados mediante métricas de clasificación y matrices de confusión, considerando especialmente el comportamiento individual de las distintas clases debido al desbalance identificado en los datos.

\begin{thebibliography}{00}
\bibitem{kaggle_unsw}
H. Bhangale,
``UNSW Complete Dataset,''
\textit{Kaggle}, 2023. [Online]. Available:
https://www.kaggle.com/datasets/harshwardhanbhangale/unsw-complete-dataset

\bibitem{enisa2025}
European Union Agency for Cybersecurity (ENISA),
``ENISA Threat Landscape 2025,''
2025. [Online]. Available:
https://www.enisa.europa.eu/publications/enisa-threat-landscape-2025

\bibitem{kaufman2012leakage}
S. Kaufman, S. Rosset, C. Perlich, and O. Stitelman,
``Leakage in data mining: Formulation, detection, and avoidance,''
\textit{ACM Transactions on Knowledge Discovery from Data (TKDD)},
vol. 6, no. 4, pp. 1--21, 2012.

\bibitem{repo}
F. A. Lanche Pacsi et al.,
``ml-project,''
GitHub, 2026. [Online]. Available:
\url{https://github.com/FabricioLanche/ml-project}

\bibitem{mehrabi2021survey}
N. Mehrabi, F. Morstatter, N. Saxena, K. Lerman, and A. Galstyan,
``A survey on bias and fairness in machine learning,''
\textit{ACM Computing Surveys (CSUR)},
vol. 54, no. 6, pp. 1--35, 2021.

\bibitem{more2024enhanced}
S. More, M. Idrissi, H. Mahmoud, and A. T. Asyhari,
``Enhanced Intrusion Detection Systems Performance with UNSW-NB15 Data Analysis,''
\textit{Algorithms},
vol. 17, no. 2, p. 64, 2024, doi: 10.3390/a17020064.

\bibitem{moustafa2015unsw}
N. Moustafa and J. Slay,
``UNSW-NB15: A comprehensive data set for network intrusion detection systems (UNSW-NB15 network data set),''
in \textit{2015 Military Communications and Information Systems Conference (MilCIS)},
Canberra, Australia, 2015, pp. 1--6.

\bibitem{moustafa2016evaluation}
N. Moustafa and J. Slay,
``The evaluation of Network Anomaly Detection Systems: Statistical analysis of the UNSW-NB15 data set and the comparison with the KDD99 data set,''
\textit{Information Security Journal: A Global Perspective},
vol. 25, no. 1--3, pp. 18--31, 2016.

\end{thebibliography}
\end{document}